





\section{Omitted Definitions  }
\begin{definition}[Left-perfect matching]
    \label{definition:left-perfect-matching}
    A matching $M$ on a bipartite graph is \emph{left-perfect} if for every vertex $v$ in the left set of the bipartite graph, there is an edge in $M$ incident to $v$.
\end{definition}

\section{Omitted proofs from Section~\ref{sec:conventional-algorithms}}

\naiveapproach*
\begin{proof}
The distortion of the algorithm follows directly from Corollary~\ref{corollary:lambda-ell-star-lower-bound}. The number of distance queries consists of three parts: First, computing the initial solution $\InitSolSet$ takes $\frac{k^2 - k}{2}$ distance queries by Theorem~\ref{theorem:ordinal-k-center-algorithm}. Second, querying all distances between the centers in $\InitSolSet$ and the groups in $\MultiGroupSet$ takes $k^2$ distance queries. Third, for any fixed $\ell$, the algorithm makes $\ell$ distance queries to compute $\cost(\InitSolSet_{\ell})$, leading to a total of $\frac{k(k+1)}{2}$ queries across all $\ell$. Summing these, the algorithm takes $2 k^2 \in \bigO(k^2)$ distance queries.
\end{proof}

\section{Omitted proofs from Section~\ref{sec:analysis-ordinal-fair-k-center-with-limited-distance-queries}}



\lambdaellstarlowerboundsmallerqueries* 
\begin{proof}
    Our proof utilizes Lemma 3.6 from Burkhardt et al.~\cite{burkhardt2024low}.
    Let us first formally restate this lemma using the notation of our paper. 
    \begin{lemma}[Lemma 3.6~\cite{burkhardt2024low}]
        \label{lemma:init-sol-mem-lower-bound-burkhardt}
        Let $\{\InitSolMem\} = \InitSolSet_{i+1} \setminus \InitSolSet_{i}$ be the center selected at the $i+1$-th step. 
        It holds that $d(\InitSolMem, \InitSolSet_{i}) \geq \frac{1}{2} \max_{u \in U} d(u, \InitSolSet_{i})$. 
    \end{lemma}

We start with proving the $\InitSolSet$ is a progressive $4$-cover for $\Ical$.
Let $\OptSolSet$ be the optimal solution for $\Ical$.
Let $\ell \in [k]$ be the largest index such that $\InitSolSet_{\ell}$ hits each part of $\mathbf{\Pi}^*$ at most once.
If $\ell = k$, then by triangle inequality, for any $u \in U$, and suppose its closest center in $\InitSolSet$ is $\InitSolMem$, and $\InitSolMem$ is optimal cluster with center $\SolMem^*$. It holds that $d(u, \InitSolSet) = d(u, \InitSolMem) \leq d(u, \SolMem^*) + d(\SolMem^*, \InitSolMem) \leq 2\cdot \cost(\OptSolSet)$.
If $\ell < k$, it indicates that $\InitSolMem_{\ell+1}$ is the first center that hits a optimal cluster 
more than once.  
Let $\InitSolMem_{i}$ where $i\leq \ell$ be a previous center that also hits the same optimal cluster. 
It holds that for any $u \in U$, it holds that $d(u, \InitSolSet_{\ell}) \leq 2 d(\InitSolMem_{\ell+1}, \InitSolSet_{\ell}) \leq 2 d(\InitSolMem_{\ell + 1}, \InitSolMem_{i}) \leq 4 \cost(\OptSolSet)$.
The first inequality holds as $\InitSolMem_{\ell+1}$ is a $\frac{1}{2}$-approximate futhest point to $\InitSolSet_{\ell}$, according to Lemma~\ref{lemma:init-sol-mem-lower-bound-burkhardt}.
The second inequality holds as $\InitSolMem_{i}$ is one of the centers in $\InitSolSet_{\ell}$. 
The third inequality holds by triangle inequality and as both $\InitSolMem_{\ell+1}$ and $\InitSolMem_{i}$ are in the same optimal cluster. 
Therefore, $\InitSolSet$ is a progressive $4$-cover for $\Ical$. 
The critical index $\ell^*$ is the $\ell$ we define above, i.e., the largest index such that $\InitSolSet_{\ell}$ hits each part of $\mathbf{\Pi}^*$ at most once.

Next, we prove that $\lambda_{\ell^*} \leq \cost(\OptSolSet)$. 
Since $\InitSolSet_{\ell}$ hits each part of $\mathbf{\Pi}^*$ at most once, it follows that each $\InitSolMem \in \InitSolSet_{\ell}$ can be uniquely mapped to each group $\MultiGroup \in \MultiGroupSet$; it implies $d(\InitSolMem, \MultiGroup) \leq \cost(\OptSolSet)$.
Let us recall the definition of $(\ell,\lambda)$-projection graph $H^{\ell}_{\lambda} = (\InitSolSet_{\ell} \cup \MultiGroupSet, E_{\lambda})$.
If there exists a left perfect matching on $H^{\ell^*}_{\lambda}$, it indicates that for any $\InitSolMem \in \InitSolSet_{\ell^*}$, there exists a $\MultiGroup \in \MultiGroupSet$ such that $d(\InitSolMem, \MultiGroup) \leq \lambda$.
This implies that when $\lambda = \cost(\OptSolSet)$, there exists a left perfect matching on $H^{\ell^*}_{\lambda}$.
Also because that $\lambda_{\ell^*}$ is the minimum possible value to guarantee the existence of a left perfect matching on $H^{\ell^*}_{\lambda_{\ell^*}}$, it follows that $\lambda_{\ell^*} \leq \cost(\OptSolSet)$.
\end{proof}

\monotonicitypropertiesoflambdaell* 
\begin{proof}
    We prove that $\lambda_{\ell}$ is a non-decreasing function in $\ell$ and $\cost(\InitSolSet_{\ell})$ is a non-increasing function in $\ell$ by the way they are defined. 
    Let $\ell < \ell'$.

    Recall $\cost(\InitSolSet_{\ell}) = \max_{u \in U} d(u, \InitSolSet_{\ell}) \geq \max_{u \in U} d(u, \InitSolSet_{\ell'}) = \cost(\InitSolSet_{\ell'})$. 
    The inequality $\geq$ holds as $\InitSolSet_{\ell} \subset \InitSolSet_{\ell'}$, therefore, $d(u, \InitSolSet_{\ell}) \geq d(u, \InitSolSet_{\ell'})$. 
    This implies that $\cost(\InitSolSet_{\ell})$ is a non-increasing function in $\ell$.

    Recall that $H^{\ell}_{\lambda} = (\InitSolSet_{\ell} \cup \MultiGroupSet, E_{\lambda})$, where $E_{\lambda} = \{(\InitSolMem, \MultiGroup) \mid d(\InitSolMem, \MultiGroup) \leq \lambda\}$. 
    $\lambda_{\ell}$ is the minimum $\lambda$ such that there exists a left perfect matching on $H^{\ell}_{\lambda}$.
    Since $\ell < \ell'$, it follows that $\InitSolSet_{\ell} \subset \InitSolSet_{\ell'}$, therefore, whenever there exists a left perfect matching on $H^{\ell'}_{\lambda_{\ell'}}$, there exists a left perfect matching on $H^{\ell}_{\lambda_{\ell'}}$.
    It implies that $\lambda_{\ell} \leq \lambda_{\ell'}$.
\end{proof}

\monotonicityoflambdacost*
\begin{proof}[Proof of Lemma~\ref{lemma:monotonicity-of-lambda-and-cost-init-sol-set}]
    Since by Lemma~\ref{lemma:monotonicity-properties-of-lambda-ell}, $\lambda_{\ell}$ is a non-de\-creas\-ing function and $\cost(\InitSolSet_{\ell})$ is a non-increasing function, w.r.t. $\ell$, we have three cases.
    
    Case (a): When $4 \lambda_{\ell} \geq \cost(\InitSolSet_{\ell})$ for all $\ell \in [k]$, it follows that $4 \lambda_{1} \geq \cost(\InitSolSet_{1})$.
    Furthermore, we have $\lambda_1 \le \lambda_{\ell^*} \le  \cost(\OptSolSet)$, due to~\Cref{lemma:lambda-ell-star-lower-bound-smaller-queries}.

    Case (b): When $4 \lambda_{\ell} \leq \cost(\InitSolSet_{\ell})$ for all $\ell \in [k]$, it follows that $4 \lambda_{k} \leq \cost(\InitSolSet_{k})$. 
    Furthermore, we have $\cost(\InitSolSet_{k}) \leq \cost(\InitSolSet_{\ell^*}) \leq 4\,\cost(\OptSolSet)$,  due to~\Cref{lemma:lambda-ell-star-lower-bound-smaller-queries}.

    Case (c): If the above two cases are false then, there exists $2 \leq \ell' \leq k-1$ such that for any $\ell \leq \ell'$, it holds that $4 \lambda_{\ell} \leq \cost(\InitSolSet_{\ell})$ and for any $\ell \geq \ell' + 1$, it holds that $4 \lambda_{\ell} \geq \cost(\InitSolSet_{\ell})$.
    Next, we show that $\min \{\cost(\InitSolSet_{\ell'}), 4 \lambda_{\ell' + 1}\} \leq 4\,\cost(\OptSolSet)$. From \Cref{lemma:lambda-ell-star-lower-bound-smaller-queries}, we have that for the critical index $\ell^*$, it holds that $\max\{4\lambda_{\ell^*}, \cost(\InitSolSet_{\ell^*})\}\leq 4\, \cost(\OptSolSet)$. First, consider the case when $\ell^* \le \ell'$, and note that since, $ \cost(\InitSolSet_{\ell^*}) \leq 4\, \cost(\OptSolSet)$, we have that $\cost(\InitSolSet_{\ell'}) \le  \cost(\InitSolSet_{\ell^*})\leq 4\, \cost(\OptSolSet)$, as desired. Now, consider the case when $\ell^* \geq \ell'+1$. Again using~\Cref{lemma:lambda-ell-star-lower-bound-smaller-queries}, we have that $4\lambda_{\ell'+1} \leq 4\lambda_{\ell'+1} \leq 4\, \cost(\OptSolSet)$, as required.
\end{proof}

\getfeasiblesolution*
\begin{proof}
    The solution is a feasible solution because $\abs{\SolSet_{\hat{\ell}}} \leq k$ and $\abs{\SolSet_{\hat{\ell}} \cap \MultiGroup_i} \geq \alpha_i$ for all $i \in [k]$.
    The distortion of the algorithm is proved by triangle inequality and Theorem~\ref{lemma:binary-search-on-ell}. 
    By Theorem~\ref{lemma:binary-search-on-ell}, we know that $\lambda_{\hat{\ell}} \leq \cost(\OptSolSet)$ and $\cost(\InitSolSet_{\hat{\ell}}) \leq 4\, \cost(\OptSolSet)$. 
    It follows that $\cost(\SolSet_{\hat{\ell}}) \leq \lambda_{\hat{\ell}} + \cost(\InitSolSet_{\hat{\ell}}) \leq 5\, \cost(\OptSolSet)$ through the left-perfect matching. 
\end{proof}

\constructgraphqueries*
\begin{proof}
    Recall that to construct the edge set $E_{\lambda}$ of $H^{\ell}_{\lambda}$, we need to determine for each center $\SolMem_{i'} \in \InitSolSet_{\ell}$ and each group $\MultiGroup_{j'} \in \MultiGroupSet$ whether $d(\SolMem_{i'}, \MultiGroup_{j'}) \leq \lambda$.
    For any center $\SolMem_{i'} \in \InitSolSet_{\ell}$, we can leverage the fact that the groups in $\MultiGroupSet$ are ordered by their distances to $\SolMem_{i'}$ in non-decreasing order (denoted by $\succ_{\SolMem_{i'}}$). 
    We perform a binary search to find the rightmost group $\MultiGroup_{j^*}$ in this ordering such that $d(\SolMem_{i'}, \MultiGroup_{j^*}) \leq \lambda$.
    This binary search requires at most $\log(k)$ distance queries.
    Once we find $\MultiGroup_{j^*}$, we can determine all edges incident to $\SolMem_{i'}$: there is an edge between $\SolMem_{i'}$ and $\MultiGroup_{j'}$ if and only if $\MultiGroup_{j'} \preceq_{\SolMem_{i'}} \MultiGroup_{j^*}$.
    This is because all groups $\MultiGroup_{j'}$ with $\MultiGroup_{j'} \preceq_{\SolMem_{i'}} \MultiGroup_{j^*}$ have distance at most $\lambda$ (by the ordering property), while all groups $\MultiGroup_{j'}$ with $\MultiGroup_{j'} \succ_{\SolMem_{i'}} \MultiGroup_{j^*}$ have distance strictly greater than $\lambda$.
    Since we perform this procedure for each of the $\ell$ centers in $\InitSolSet_{\ell}$, the total number of distance queries is at most $\ell \log(k)$.
\end{proof}

\evaluatepredicatequeries*
\textsc{Proof of~\Cref{lemma:evaluate-predicate-queries}.}
   For any fixed $\lambda$, it takes at most $\ell \log(k)$ distance queries to construct the bipartite graph $H^{\ell}_{\lambda}$ according to Lemma~\ref{lemma:construct-graph-queries}.
    In addition, $\cost(\InitSolSet_{\ell})$ takes $\ell$ distance queries to compute.
    Checking whether there exists a left-perfect matching on $H^{\ell}_{\lambda}$ does not take additional distance queries. 
\hfill\qed


\reducesearchspace*
\begin{proof}
    The proof closely follows the reasoning in Lemma~\ref{lemma:construct-graph-queries}, with only a slight variation in the binary search condition.  
\end{proof}

\medianmedian* 
\textsc{Proof.}
First, we make an observation on $\pidx$ that $\sum_{i=\pidx}^{k} \abs{D_{\pi_{i}}^{(t)}} \geq \frac{L^{(t)}}{2}$ also holds. 
This is because that $\pidx$ is the smallest $j$ such that $\sum_{i=1}^{j} \abs{D_{\pi_{i}}^{(t)}} \geq \frac{L^{(t)}}{2}$, it follows that $\sum_{i=1}^{\pidx-1} \abs{D_{\pi_{i}}^{(t)}} < \frac{L^{(t)}}{2}$.
Thus $\sum_{i=\pidx}^{k} \abs{D_{\pi_{i}}^{(t)}} = L^{(t)} - \sum_{i=1}^{\pidx-1} \abs{D_{\pi_{i}}^{(t)}} > \frac{L^{(t)}}{2}$.
Next, we make a case distinction and let $\tau = \MED(D_{\pi_{\pidx}}^{(t)})$. 

Case 1: $H^{\ell}_{\tau}$ has a left-perfect matching.
The algorithm records $\tau$ as the current best feasible result. 
All the distance values in $D^{(t)}$ that are equal or greater than $\tau$ are thus removed from the search space. 
Recall that $D_{\pi_i}^{(t)}$ are sorted in non-decreasing order according to the distance value of $\tau$. 
At least $\frac{1}{2} \sum_{i=\pidx}^{k} \abs{D_{\pi_{i}}^{(t)}} \geq \frac{L^{(t)}}{4}$ distance values are removed from the search space.


Case 2: $H^{\ell}_{\tau}$ has no left-perfect matching. 
All the distance values in $D^{(t)}$ that are less than or equal to $\tau$ are removed from the search space.
Thus at least $\frac{1}{2} \sum_{i=1}^{\pidx} \abs{D_{\pi_{i}}^{(t)}} \geq \frac{L^{(t)}}{4}$ distance values are removed from the search space.

Thus, the size of the search space is reduced by at least $\frac{1}{4}$.
\hfill\qed



    
















        
